\section{Methodology}
\label{sec:method}

\subsection{Problem formulation}

We study closed-book short-answer QA.
A model receives a question $q$ and generates an answer $a$ with no supporting passage.
We label the answer $y=1$ (\emph{hallucinated}) if it is factually wrong and $y=0$ (\emph{factual}) if it is right, and we drop abstentions.
Our question is whether the attention entropy the model shows while producing $a$ predicts $y$.

\para{Hypotheses.} We fixed three claims before running any model:
\begin{itemize}[leftmargin=*,itemsep=0pt,topsep=0pt]
    \item \Hone: mean attention entropy over the final three layers is 15--30\% higher for hallucinated than for factual answers. We count it as supported only if the relative difference and its 95\% CI overlap that band with the right sign.
    \item \Htwo: logistic regression on layer-wise entropy alone reaches a cross-validated AUROC above 0.75, pooled over domains. We count it as supported if the lower bound of the 95\% CI exceeds 0.75.
    \item \Hthree: within-domain AUROC is highest for biography and lowest for science.
\end{itemize}

\subsection{Models and generation}

Attention maps are required, so we use two local open-weight models: \texttt{Qwen2.5-7B-Instruct} (28 layers $\times$ 28 heads; \qwen) and \texttt{Llama-3.1-8B-Instruct} (32 layers $\times$ 32 heads; \llama).
We run both in float32 with eager attention.
We generate with greedy decoding and at most 24 new tokens, using each model's chat template and a fixed system prompt that asks for a short answer with no explanation and a best guess (full text in \secref{app:details}).

\subsection{Data and labels}

We sample 1,500 questions per domain, 7,500 in total, with a fixed seed (\tabref{tab:data}).
Biography, geography, and creative works come from \popqa relations \citep{mallen2023popqa}, science from \sciq \citep{welbl2017sciq}, and trivia from \triviaqa \citep{joshi2017triviaqa}.

\begin{table}[t]
    \centering
    \caption{Data. Each domain has 1,500 questions (seed 42). The hallucination rate is the share of non-abstaining answers labeled wrong; it ranges from 23\% to 94\% across domains, which is the source of the domain confound in \secref{sec:h2}.}
    \label{tab:data}
    \resizebox{\textwidth}{!}{%
    \begin{tabular}{@{}llcc@{}}
        \toprule
        & & \multicolumn{2}{c}{\textbf{Hallucination rate}} \\
        \cmidrule(lr){3-4}
        \textbf{Domain} & \textbf{Source} & \qwen & \llama \\
        \midrule
        biography & \popqa: occupation, place of birth, father, mother & 0.77 & 0.64 \\
        geography & \popqa: country, capital, capital of & 0.38 & 0.33 \\
        creative works & \popqa: author, director, producer, screenwriter, composer & 0.94 & 0.80 \\
        science & \sciq (train split, question only) & 0.23 & 0.27 \\
        trivia & \triviaqa (rc.nocontext validation) & 0.42 & 0.25 \\
        \bottomrule
    \end{tabular}
    }
\end{table}


\para{Labels.}
We label an answer factual if a gold alias string-matches the normalized answer or a reference-based LLM judge grades it correct; otherwise we label it hallucinated.
The judge is \judge with thinking disabled and greedy decoding; it sees the question, the gold answers, and the model answer.
We drop abstentions, which we identify with a regular-expression rule combined with the judge's verdict.
This removes 30 examples for \qwen and 50 for \llama, leaving 7,470 and 7,450.

The judge and string match agree on 92.5\% (\qwen) and 92.1\% (\llama) of examples.
Agreement is lowest in science (76\% and 78\%), where gold answers are free-form phrases.
In a manual audit of 60 randomly chosen primary labels we found 2 clear errors and 3 debatable cases with an ambiguous question or gold answer.
We repeat every headline number under judge-only and string-match-only labels.

\subsection{Attention-entropy features}

We run one teacher-forced forward pass over the prompt followed by the generated answer.
Let $\alpha^{(\ell,h)}_{p,j}$ be the attention weight that head $h$ in layer $\ell$ places from position $p$ on position $j \le p$.
For every position $p$ that produces an answer token, the attention entropy in nats is
\begin{equation}
    H^{(\ell,h)}_p = -\sum_{j \le p} \alpha^{(\ell,h)}_{p,j} \log \alpha^{(\ell,h)}_{p,j}.
    \label{eq:entropy}
\end{equation}
We aggregate over positions in two ways.
The primary feature, \emph{answer mean}, averages $H^{(\ell,h)}_p$ over all answer-producing positions.
The \emph{first-token} feature uses only the position that produces the first answer token; it depends on the prompt alone, so it cannot leak answer length.
A layer feature is the mean over heads, which gives $L=28$ or $32$ features per example.
The \Hone statistic is the mean of the last three layer features.
Because $H_p \le \log n_p$ for $n_p$ attended positions, we also compute a length-normalized variant $H_p / \log n_p$.

\subsection{Baselines and controls}

We compare layer-wise entropy against the following feature sets, all fed to the same classifier:
\begin{itemize}[leftmargin=*,itemsep=0pt,topsep=0pt]
    \item \textbf{Domain indicator only}: a one-hot encoding of the five domains. It measures how much a detector gains just from knowing which domain a question comes from.
    \item \textbf{Length only}: log prompt tokens and log answer tokens.
    \item \textbf{Logit baseline}: six features from the same forward pass, namely the mean, maximum, and first-token values of token negative log-likelihood and of next-token predictive entropy.
    \item \textbf{Combinations}: entropy + length, and entropy + logit.
    \item \textbf{Per-head entropy}: all $L \times H$ head-level entropies (784 or 1,024 features), as an upper bound for entropy-only detectors.
\end{itemize}

\subsection{Statistics}

\para{Effect sizes.}
For \Hone we report the relative difference of means, $(\bar{H}_{\text{halluc}} - \bar{H}_{\text{factual}})/\bar{H}_{\text{factual}}$, with a percentile bootstrap CI from 10,000 resamples, together with Cohen's $d$.
Per layer we run Mann--Whitney $U$ tests with Benjamini--Hochberg correction across layers \citep{benjamini1995fdr}.
To adjust for confounds we fit an OLS regression of final-three-layer entropy on the hallucination label with domain fixed effects and log prompt and answer length as covariates, using HC3 standard errors.

\para{Classifiers.}
Every classifier standardizes its features and fits an L2-regularized logistic regression with $C=1$.
We use 5-fold stratified cross-validation repeated 5 times, with standardization fit inside each fold, and compute AUROC on out-of-fold scores with a 2,000-resample bootstrap CI.
For \Hthree we train and cross-validate a separate classifier inside each domain, and we also train on one domain and test on another.
All intervals are 95\%.

\para{Compute.}
We ran all experiments on four RTX A6000 GPUs (48 GB) shared with other jobs.
Feature extraction took 24 minutes for \qwen and 28 minutes for \llama at batch size 16, judging took about 30 minutes per model, and each full analysis took about 2 minutes on CPU.
